\documentclass[12pt]{article}
\usepackage[utf8]{inputenc}
\usepackage[T1]{fontenc}
\usepackage{hyperref}
\title{Shabak Language Resources for K-12 Curriculum Development}
\author{Compiled by Claude Sonnet 5.5 (claude-sonnet-5-5) for Andr\'as Kornai}
\date{}
\begin{document}
\maketitle

\section{Introduction}
Shabak (Shabaki; ISO 639-3 sdb; Glottocode shab1251) is a Northwestern Iranian language of the Zaza--Gorani group, spoken by the Shabak community in villages east of Mosul on the Nineveh Plain of Iraq. Estimates of speakers range from 10,000--20,000 (a 1989 figure) to about 250,000 (more recent). The community was displaced by ISIS in 2014, which has heightened fears for the language. It is written, where written at all, in Arabic script. Searches found no school materials; the resources below are reference and background only.

\section{Primary Educational Resources}
No primer, textbook or school programme was found in search results. According to Minority Rights Group, the Shabak language is not taught in schools because the Kurdistan Regional Government does not recognise the Shabak as a distinct group:
\url{https://minorityrights.org/minorities/shabak/}

\section{Linguistic and Reference Materials}
\begin{itemize}
\item Wikipedia article on the Shabaki language: \url{https://en.wikipedia.org/wiki/Shabaki_language}
\item Omniglot page, which describes an Arabic-based alphabet of 28 letters: \url{https://www.omniglot.com/writing/shabaki.php}
\item Glottolog entry for sdb (endangerment status ``shifting''): \url{https://glottolog.org/resource/languoid/iso/sdb}
\end{itemize}
No grammar or dictionary was found in the search results.

\section{Religious and Community Resources}
\begin{itemize}
\item Minority Rights Group profile of the Shabak: \url{https://minorityrights.org/minorities/shabak/}
\item Refworld copy of the World Directory of Minorities entry: \url{https://www.refworld.org/reference/countryrep/mrgi/2017/en/64781}
\item Wikipedia article on the Shabaks: \url{https://en.wikipedia.org/wiki/Shabaks}
\item Joshua Project profile: \url{https://joshuaproject.net/people_groups/19500}
\end{itemize}

\section{Digital and Audio Resources}
No audio, video, corpus or software resources specific to Shabaki were found.

\section{Curriculum Development Context}
There is essentially no documented educational infrastructure for Shabak. A curriculum would have to begin with documentation: community-led standardisation of an Arabic-script orthography, recording of oral traditions, and collaboration with Shabak community representatives, who dispute classification as Kurdish. Political recognition of the Shabak as a distinct group is a precondition for school use, per the sources above.
\end{document}